% !TEX program = xelatex
% Generated by docs/latex/build.py from the sources pinned in sources.json.
\documentclass[11pt,a4paper]{article}
\usepackage[a4paper,margin=25mm]{geometry}
\usepackage{fontspec}
\setmainfont{lmroman10-regular.otf}[BoldFont=lmroman10-bold.otf,ItalicFont=lmroman10-italic.otf,BoldItalicFont=lmroman10-bolditalic.otf]
\setsansfont{lmsans10-regular.otf}[BoldFont=lmsans10-bold.otf,ItalicFont=lmsans10-oblique.otf]
\setmonofont{JuliaMono-Regular.ttf}[Path=../fonts/,Scale=0.83]
\usepackage{amsmath,amssymb,mathtools}
\usepackage{graphicx,calc,longtable,booktabs,array}
\usepackage{fvextra}
\usepackage{xurl}
\usepackage[colorlinks=true,linkcolor=blue,urlcolor=blue,citecolor=blue]{hyperref}
\providecommand{\tightlist}{\setlength{\itemsep}{0pt}\setlength{\parskip}{0pt}}
\setlength{\parindent}{0pt}
\setlength{\parskip}{6pt}
\setlength{\emergencystretch}{6em}
\begin{document}
\section{P3 C2 --- Triangular maximum depths}\label{p3-c2-triangular-maximum-depths}

Let \(B\) be the Bulgarian-solitaire move on partitions. Write \(T_k=k(k+1)/2\), and let \(D_B(n)\) be the maximum, over partitions of \(n\), of the number of moves before first reaching a cyclic partition. For every \(k\ge1\),

\[
\boxed{D_B(T_k)=k(k-1).}
\]

The upper proof covers every starting partition; the lower proof exhibits a partition attaining the bound. At triangular size the only cyclic partition is the fixed staircase \((k,k-1,\ldots,1)\). The general upper and lower bounds below are written mathematical proofs. The associated Lean development establishes triangular convergence and uniqueness for all \(k\), exact maximum depths for \(k=1,\ldots,6\), and the abstract arithmetic used in the upper proof; the pile-lifetime application has not been formalized in Lean.

\subsection{1. Sharp upper bound}\label{sharp-upper-bound}

For every \(k\ge1\) and every partition \(\lambda\) of \(T_k=k(k+1)/2\), Bulgarian solitaire reaches the staircase \((k,k-1,\ldots,1)\) in at most \(k(k-1)\) moves. Together with the explicit \href{https://github.com/mpelteshki/climbing-to-the-frontier/blob/212bba53b1818c08c40b8c4b48f636651eaf6410/cagent/p3/c2/lower-bound-proof.md}{lower-bound witness}, this proves \(D_B(T_k)=k(k-1)\) as a \textbf{written mathematical proof}. The general upper bound in this document has not been formalized in Lean; the Lean package establishes the triangular destination and finite cases.

The argument reconstructs the upper-bound half of \href{https://sc.edu/study/colleges_schools/artsandsciences/mathematics/research/imi/research/documents/1998/1998_12.pdf}{Griggs and Ho, Theorem 3.7}. It proves the needed timing facts directly using pile lifetimes. These correspond to their Proposition 3.2(1),(3) and Lemmas 3.3--3.6; their Proposition 3.2(2) is not needed for this route. We do not use Theorem 3.7 as a premise.

\subsubsection{Pile lifetimes and a sequence bound}\label{pile-lifetimes-and-a-sequence-bound}

Write \(c_i\) for the number of piles in \(B^{i-1}(\lambda)\), with \(i\ge1\). An original pile of size \(a\) exists at times \(1,\ldots,a\). The pile created by move \(i\) has size \(c_i\) and exists at times

\[
J_i=[i+1,i+c_i].
\]

Its last time is \(e_i=i+c_i\). At any time \(m\), precisely \(c_m\) of these original-pile and created-pile intervals contain \(m\). One new interval starts between times \(i\) and \(i+1\); if \(d_i\) intervals end at \(i\), then

\[
c_{i+1}=c_i+1-d_i.
\]

\begin{enumerate}
\def\labelenumi{(\arabic{enumi})}
\tightlist
\item
\end{enumerate}

In particular \(c_{i+1}\le c_i+1\); a rise by one means \textbf{no} pile dies, and a constant pair means \textbf{exactly one} pile dies.

We use a \emph{sandwich pattern} of level \(x\) and width \(L=q-p\ge2\):

\[
(c_p,c_{p+1},\ldots,c_q)=(x-1,x,\ldots,x,x+1).
\]

\begin{enumerate}
\def\labelenumi{(\arabic{enumi})}
\setcounter{enumi}{1}
\tightlist
\item
\end{enumerate}

The two rises occur at its ends. The following elementary sandwich rule will locate one. If \(i<j\), \(c_i\le x-1\), and \(c_j\ge x+1\), choose \(p\) as the last index before \(j\) with \(c_p\le x-1\), and \(q\) as the first subsequent index with \(c_q\ge x+1\). Equation (1) forces \(c_p=x-1\), \(c_{p+1}=\cdots=c_{q-1}=x\), and \(c_q=x+1\). Thus (2) occurs within \([i,j]\).

\subsubsection{The retreat lemma}\label{the-retreat-lemma}

Suppose (2) occurs and \(p>x\). Among the \(x\) recent created piles \(J_{p-x},\ldots,J_{p-1}\), at least one is absent at time \(p\), since \(c_p=x-1\). The last, \(J_{p-1}\), is present. Let \(u\in[p-x,p-2]\) be the largest index whose interval is absent. All of \(J_{u+1},\ldots,J_{p-1}\) are present at \(p\). Since \(c_{p+1}=c_p+1\), none dies at \(p\), so they remain present at \(p+1\).

Put \(h=p-u-1\) and \(y=h+1=p-u\le x\). The absence of \(J_u\) at \(p\) gives \(c_u\le h\). The presence of \(J_{u+1}\) at \(p+1\) gives \(c_{u+1}\ge h+1\). Applying \(c_{u+1}\le c_u+1\) shows

\[
c_u=h=y-1,\qquad c_{u+1}=h+1=y.
\]

\begin{enumerate}
\def\labelenumi{(\arabic{enumi})}
\setcounter{enumi}{2}
\tightlist
\item
\end{enumerate}

Moreover \(J_{u+1}\) ends at time \(u+1+y=p+1\).

If \(L=2\), the pattern also has \(c_{p+2}=c_{p+1}+1\), so no interval ends at \(p+1\), a contradiction. Hence every width-two pattern satisfies

\[
p\le x.
\]

\begin{enumerate}
\def\labelenumi{(\arabic{enumi})}
\setcounter{enumi}{3}
\tightlist
\item
\end{enumerate}

Now let \(L\ge3\). As long as no later value \(c_{u+s}\) has risen from \(y\) to \(y+1\), the consecutive intervals

\[
J_{u+s},J_{u+s+1},\ldots,J_{p+s-1}
\]

are all present at time \(p+s\), for \(1\le s\le L-1\). This holds for \(s=1\) by the choice of \(u\) and the lack of a death at \(p\). If \(c_{u+s}=y\), then \(J_{u+s}\) ends at \((u+s)+y=p+s\). For \(s\le L-2\), the pair \(c_{p+s}=c_{p+s+1}=x\) permits exactly one death. It must be this interval; all the others survive, and \(J_{p+s}\) is born. This proves the induction. At each stage the active \(J_{u+s}\) implies \(c_{u+s}\ge y\), while (1) and the preceding value \(y\) imply \(c_{u+s}\le y+1\). Hence it either stays at \(y\) or supplies the desired rise.

At \(s=L-1\), the transition from \(c_{q-1}=x\) to \(c_q=x+1\) permits no death. Thus \(J_{u+L-1}\) cannot have size \(y\), which would make it end at \(q-1\). A first rise occurs at some \(v\) with

\[
u+2\le v\le u+L-1,
\qquad(c_u,\ldots,c_v)=(y-1,y,\ldots,y,y+1).
\]

\begin{enumerate}
\def\labelenumi{(\arabic{enumi})}
\setcounter{enumi}{4}
\tightlist
\item
\end{enumerate}

It occurs by \(p+1\): otherwise both \(c_p\) and \(c_{p+1}\) would equal \(y\), contradicting their values \(x-1,x\). Thus \(u\ge p-x\), \(v\le p+1\), \(y\le x\), and the new width \(v-u<L\). This is the retreat step.

Each retreat decreases the positive integer width. It therefore stops either at a pattern with start at most its level or at width two, where (4) gives the same bound. Working backward through \(p\le u+x\), while levels never increase, gives the explicit bound

\[
\boxed{p\le x(L-1)}
\]

\begin{enumerate}
\def\labelenumi{(\arabic{enumi})}
\setcounter{enumi}{5}
\tightlist
\item
\end{enumerate}

for every sandwich pattern (2). More formally, if \(p\le x\), then \(p\le x(L-1)\) since \(L\ge2\). Otherwise induction on \(L\) gives \(p\le u+x\le y(L'-1)+x\le x(L-2)+x=x(L-1)\) for the smaller width \(L'<L\). This abstract induction is also kernel-checked as \href{https://github.com/mpelteshki/climbing-to-the-frontier/blob/212bba53b1818c08c40b8c4b48f636651eaf6410/cagent/p3/lean/ProofPursuit/P3/C2SandwichBound.lean}{\texttt{sandwich\_start\_bound}}; the lifetime argument above proves its application hypotheses in writing.

\subsubsection{A special full-width pattern}\label{a-special-full-width-pattern}

We also need one bound that uses the number \(n\) of cards. Suppose

\[
(c_p,\ldots,c_{p+k})=(k-2,k-1,\ldots,k-1,k).
\]

\begin{enumerate}
\def\labelenumi{(\arabic{enumi})}
\setcounter{enumi}{6}
\tightlist
\item
\end{enumerate}

Here \(k\ge3\). The interval \(J_p\) ends at \(p+k-2\). Since \(c_{p+k-2}=c_{p+k-1}=k-1\), it is the \textbf{only} interval ending then. No interval ends at \(p+k-1\), since the next count rises from \(k-1\) to \(k\).

At time \(p+k\), the \(k-1\) intervals \(J_{p+1},\ldots,J_{p+k-1}\) are all present. There is exactly one other pile. If it is original, its original size is at least \(p+k\), and hence \(n\ge p+k\). If it is \(J_1\), then \(1+c_1\ge p+k\), and \(c_1\le n\) gives \(n+1\ge p+k\).

The only other possibility would be \(J_i\) for some \(2\le i\le p-1\); \(J_p\) has already died. Such a \(J_i\) is also the unique extra pile at time \(p+k-1\), alongside \(J_{p+1},\ldots,J_{p+k-2}\). Hence \(J_{i-1}\) is absent at that time, so \((i-1)+c_{i-1}\le p+k-2\). But \(J_i\) survives to \(p+k\), and (1) yields

\[
(i-1)+c_{i-1}\ge(i-1)+(c_i-1)\ge p+k-2.
\]

Equality follows, making \(J_{i-1}\) a second interval ending at \(p+k-2\), contrary to the uniqueness of \(J_p\). We have proved

\[
\boxed{p+k\le n+1}.
\]

\begin{enumerate}
\def\labelenumi{(\arabic{enumi})}
\setcounter{enumi}{7}
\tightlist
\item
\end{enumerate}

\subsubsection{Finding the final pattern}\label{finding-the-final-pattern}

Every trajectory at triangular size eventually reaches the staircase, by \href{https://github.com/mpelteshki/climbing-to-the-frontier/blob/212bba53b1818c08c40b8c4b48f636651eaf6410/cagent/p3/lean/ProofPursuit/P3/TriangularGeneral.lean}{triangular convergence}. Let \(t\) be the \textbf{first} move at which it does, so \(B^t(\lambda)\) is the staircase. If \(t=0\), there is nothing to prove. Otherwise \(c_{t+1}=k\). The fresh pile made on move \(t\) has size \(c_t\), so \(c_t\le k\). If \(c_t=k\), then its predecessor has \(k\) piles and must itself be the staircase: after removing one from each, its surviving piles must be exactly \(k-1,k-2,\ldots,1\), with one exhausted pile. This contradicts minimality of \(t\). Equation (1) now forces \(c_t=k-1\), and the rise to \(c_{t+1}=k\) entails no death at time \(t\).

Assume \(t\ge k+1\). Of the \(k\) intervals \(J_{t-k},\ldots,J_{t-1}\), at least one is absent at time \(t\), because only \(k-1\) piles exist. The last, \(J_{t-1}\), is present. Choose the largest absent index \(u\in[t-k,t-2]\). Repeating the argument leading to (3), now with \(p=t\) and \(x=k\), gives

\[
c_u=t-u-1,\qquad c_{u+1}=t-u.
\]

\begin{enumerate}
\def\labelenumi{(\arabic{enumi})}
\setcounter{enumi}{8}
\tightlist
\item
\end{enumerate}

If \(u\ge t-k+1\), then \(c_u\le k-2\). Apply the sandwich rule between \(u\) and \(t+1\) at level \(k-1\). It produces a type-II pattern \(k-2,k-1,\ldots,k-1,k\) with endpoints \(p\ge t-k+1\), \(q\le t+1\), and width \(q-p\le k\). When that width is exactly \(k\), it is the special pattern (7).

The remaining case is \(u=t-k\), giving \(c_{t-k}=k-1\) and \(c_{t-k+1}=k\). If some \(c_j\le k-2\) for \(t-k+2\le j\le t-1\), the same sandwich rule again produces a type-II pattern within \([t-k+1,t+1]\). If instead some such \(c_j\ge k+1\), the rule at level \(k\), between \(t-k\) and \(j\), produces a pattern

\[
(c_p,\ldots,c_q)=(k-1,k,\ldots,k,k+1),
\quad t-k\le p<q\le t-1,
\quad q-p\le k-1.
\]

\begin{enumerate}
\def\labelenumi{(\arabic{enumi})}
\setcounter{enumi}{9}
\tightlist
\item
\end{enumerate}

Finally suppose every intermediate \(c_j\) is \(k-1\) or \(k\). By the choice \(u=t-k\), all \(J_{t-k+1},\ldots,J_{t-1}\) are present at time \(t\). They account for all \(k-1\) piles, and all survive to \(t+1\); the only new pile is \(J_t\), of size \(k-1\). At time \(t+1\), each older \(J_j\) has remaining size \(j+c_j-t\le (t-1)+k-t=k-1\). No pile can have size \(k\), contradicting the staircase. Thus either (10) occurs, or a type-II pattern occurs with width at most \(k\).

\subsubsection{The deadline}\label{the-deadline}

For \(k=1\), the only partition is fixed. For \(k=2\), the three partitions of \(3\) have depths \(1,0,2\), respectively: \((3)\to(2,1)\), \((2,1)\) is fixed, and \((1,1,1)\to(3)\to(2,1)\). Now let \(k\ge3\). If \(t\le k\), then \(t\le k(k-1)\).

Suppose \(t\ge k+1\). If a type-II pattern has full width \(k\), it is (7), and its allowed index range forces \(p=t-k+1\), \(q=t+1\). By (8), \(t=p+k-1\le n=T_k\le k(k-1)\), the last inequality holding for \(k\ge3\).

For every other pattern just found, its level \(x\) is \(k\) or \(k-1\), its width \(L\) is at most \(k-1\), and its start satisfies \(p\ge t-k\). Equation (6) gives

\[
t\le p+k\le x(L-1)+k\le k(k-2)+k=k(k-1).
\]

The lower-bound witness attains this deadline, so the maximum first cyclic depth at \(T_k\) is exactly \(k(k-1)\).

\subsection{2. Attaining construction}\label{attaining-construction}

Let \(k\ge2\), \(n=T_k=k(k+1)/2\), and \(B\) be the Bulgarian-solitaire move. The partition

\[
\lambda^{(0)}=(k-1,k-1,k-2,\ldots,2,1,1)
\]

has first cyclic depth exactly \(k(k-1)\). The notation includes three copies of \(1\) when \(k=2\), so \(\lambda^{(0)}=(1,1,1)\) in that case. Consequently the maximum depth satisfies \(D_B(T_k)\ge k(k-1)\). For \(k=1\), the sole partition \((1)\) is fixed and gives the same lower bound \(0\). This proves the lower-bound half only; it does not establish an upper bound for arbitrary starting partitions.

Use zero-based columns \(j\) and one-based rows \(i\). A cell has diagonal index \(i+j\). Let \(S_k\) be the staircase diagram: every cell on diagonals \(1,\ldots,k\) is occupied, so column \(j<k\) has height \(k-j\). For \(0\le p<k\) and \(0\le q\le k\), form \(D(p,q)\) by deleting the diagonal-\(k\) cell in column \(p\) and adding the diagonal-\((k+1)\) cell in column \(q\). It has \(T_k\) cells. The added cell leaves a vertical gap when \(q=p\). Otherwise its column heights, including a possible zero at the end, are

\[
h_j=k-j-\mathbf1_{j=p}+\mathbf1_{j=q}\qquad(0\le j\le k).
\]

These heights are weakly decreasing exactly when \(q\ne p+1\). Indeed the only way the deletion and addition can reverse an adjacent pair is to decrease column \(p\) and increase column \(p+1\). Thus \(D(p,q)\) is a Ferrers partition whenever \(q\notin\{p,p+1\}\). In particular, \(D(0,k)=\lambda^{(0)}\).

Before sorting, a Bulgarian move maps a card \((i,j)\) with \(i>1\) to \((i-1,j+1)\), and a top card \((1,j)\) to \((j+1,0)\). Hence on diagonal \(w\) it rotates column \(j\) to \(j+1\pmod w\). Every full diagonal of \(D(p,q)\) remains full. The hole and added card move respectively to columns \(p+1\pmod k\) and \(q+1\pmod{k+1}\). If the resulting \(D\) is Ferrers, it is already in decreasing pile order, so the sorting part of \(B\) changes nothing.

Put \(T=k(k-1)\) and, for \(0\le t<T\), set

\[
p_t=t\bmod k,\qquad q_t=(k+t)\bmod(k+1).
\]

Write \(t=ak+b\), where \(0\le a\le k-2\) and \(0\le b<k\). Then \(p_t=b\) and, because \(k\equiv-1\pmod{k+1}\),

\[
q_t=\begin{cases}
b-a-1,&b\ge a+1,\\
k+b-a,&b\le a.
\end{cases}
\]

In the first case \(q_t<b\); in the second \(q_t-b=k-a\ge2\). Therefore \(q_t\ne p_t\) and \(q_t\ne p_t+1\) throughout \(0\le t<T\). Each \(D(p_t,q_t)\) is thus a partition. Starting with \(D(p_0,q_0)=\lambda^{(0)}\), the diagonal rotation and the Ferrers criterion give, by induction,

\[
B^t(\lambda^{(0)})=D(p_t,q_t)\qquad(0\le t<T).
\]

At \(t=T-1=(k-2)k+(k-1)\), we have \(p_t=k-1\) and \(q_t=0\), so the partition is \((k+1,k-1,k-2,\ldots,2)\), with the descending tail empty for \(k=2\). Its next \emph{unsorted} move is \((k-1,k,k-2,\ldots,1)\): this is \(D(0,1)\), since \(p_T=0\) and \(q_T=1\). Sorting its first two entries yields the staircase \((k,k-1,\ldots,1)\), which is fixed by \(B\). This also shows explicitly that the first sorting change occurs on the move numbered \(T\).

For every \(t<T\), \(D(p_t,q_t)\) still has an occupied cell on diagonal \(k+1\), whereas the staircase has none. The staircase is the \textbf{only cyclic partition} of \(T_k\), by the general triangular uniqueness theorem proved in \href{https://github.com/mpelteshki/climbing-to-the-frontier/blob/212bba53b1818c08c40b8c4b48f636651eaf6410/cagent/p3/lean/ProofPursuit/P3/TriangularGeneral.lean}{TriangularGeneral.lean} and explained in the \href{https://github.com/mpelteshki/climbing-to-the-frontier/blob/212bba53b1818c08c40b8c4b48f636651eaf6410/cagent/p3/c1/written-proof.md}{C1 proof}. Thus no earlier state is cyclic, and the first cyclic depth of this witness is exactly \(T=k(k-1)\).

The witness and lower-bound result appear in \href{https://sc.edu/study/colleges_schools/artsandsciences/mathematics/research/imi/research/documents/1998/1998_12.pdf}{Griggs and Ho, Theorem 3.1}. The modular calculation above supplies an explicit check that the unsorted diagonal rotations remain valid partitions at every earlier step.

\subsection{Verification and provenance}\label{verification-and-provenance}

The general triangular destination and cyclic uniqueness are proved in \href{https://github.com/mpelteshki/climbing-to-the-frontier/blob/212bba53b1818c08c40b8c4b48f636651eaf6410/cagent/p3/lean/ProofPursuit/P3/TriangularGeneral.lean}{TriangularGeneral.lean}. The exact finite maximum-depth theorems are in the \href{https://github.com/mpelteshki/climbing-to-the-frontier/tree/212bba53b1818c08c40b8c4b48f636651eaf6410/cagent/p3/lean/ProofPursuit/P3}{P3 Lean package}, and the abstract sandwich induction is in \href{https://github.com/mpelteshki/climbing-to-the-frontier/blob/212bba53b1818c08c40b8c4b48f636651eaf6410/cagent/p3/lean/ProofPursuit/P3/C2SandwichBound.lean}{C2SandwichBound.lean}. The \href{https://github.com/mpelteshki/climbing-to-the-frontier/blob/212bba53b1818c08c40b8c4b48f636651eaf6410/cagent/p3/evidence/verification.json}{recorded Lean replay} is finite or abstract evidence as stated, not a Lean formalization of the all-\(k\) depth result.

The sharp result is classical. The proof cites Griggs and Ho's original preprint for historical context and reconstructs its needed upper-bound steps rather than assuming the target theorem.
\end{document}
